\section*{РЕФЕРАТ}
\addcontentsline{toc}{section}{РЕФЕРАТ}

Работа содержит \pageref{lastpage} страниц, \ldots{} рисунков, \ldots{} таблиц,
\ldots{} источников, \ldots{} приложений.

\textbf{Ключевые слова:} PTX ISA, CUDA, программный эмулятор,
LD\_PRELOAD, синтаксический анализ, SIMT, дивергенция потоков, PDOM-реконвергенция,
профилировщик, глобальная когерентность памяти.

Данная выпускная квалификационная работа посвящена разработке программного
эмулятора инструкций PTX (Parallel Thread Execution) --- промежуточного
представления CUDA-программ компании NVIDIA. Цель работы — создать инструментальное средство, обеспечивающее прозрачный перехват CUDA Runtime
API посредством механизма \texttt{LD\_PRELOAD}, синтаксический анализ PTX-текста,
исполнение набора арифметических, логических инструкций и инструкций работы
с памятью, а также управление дивергенцией потоков на основе алгоритма
PDOM-реконвергенции и сбор метрик производительности на уровне инструкций.

Механизм перехвата реализован в виде динамической библиотеки
\texttt{libemuruntime.so}, подставляемой вместо \texttt{libcudart.so}
без модификации целевого бинарного файла. Архитектура эмулятора
основана на таблице диспетчеризации с диспетчеризацией по имени мнемоники
и формальной операционной семантикой. Механизм дивергенции использует стек фреймов
реконвергенции на основе метода постдоминаторов (PDOM).

Система профилирования собирает пять метрик (\texttt{branch\_efficiency},
\texttt{bank\_conflicts}, \texttt{global\_mem\_transactions},
\texttt{global\_coalescing}, \texttt{write\_ops}) через кодогенерируемые хуки, определяемые
JSON-манифестом, и формирует самодостаточный HTML-отчёт с аннотацией
строк исходного кода.

Корректность эмулятора подтверждена прохождением семи интеграционных
тестов на реальных CUDA-ядрах (\texttt{vadd}, \texttt{gelu},
\texttt{softmax}, \texttt{mmul}, \texttt{attention}, \texttt{fft},
\texttt{conjugate\_gradient}); отклонение результатов от CPU-эталона
не превышает $10^{-3}$. Корректность пяти метрик профилировщика подтверждена на тестовых ядрах с аналитически вычисленными ожидаемыми значениями.

\newpage
